%------------------------------ Ejercicio 1.3.A -------------------------------
\begin{nejer}
Muestra que cualesquiera dos objetos iniciales son isomorfos. Muestra que cualesquiera dos objetos finales son isomorfos
\end{nejer}
\begin{sol}
Sean $A,B$ objetos iniciales. Existen únicos $f\colon A\to B$ y $g\colon B\to A$. Como $g\circ f$ y $\id_A$ son morfismos en $\hom(A,A)$, la unicidad implica $g\circ f=\id_A$. Asimismo, $f\circ g$ y $\id_B$ son morfismos en $\hom(B,B)$, así que la unicidad implica $f\circ g=\id_B$. Luego $f  \colon A \to B$ es isomorfismo.

Sean $A,B$ objetos finales. Existen únicos $f\colon A\to B$ y $g\colon B\to A$. Como $g\circ f$ y $\id_A$ son morfismos en $\hom(A,A)$, la unicidad implica $g\circ f=\id_A$. Asimismo, $f\circ g$ y $\id_B$ son morfismos en $\hom(B,B)$, así que la unicidad implica $f\circ g=\id_B$. Luego $f \colon A \to B$ es isomorfismo.
\end{sol}
